%%%%%%%%%%%%%%%%%%%%%%%%%%%%%%%%%%%%%%%%%%%%%%%%%%%%%%%%%%%%
\section{Conclusion}
\label{sec:conclusion}

We presented authenticated workflows—the first complete trust layer for enterprise agentic AI. Our key architectural insight: positioning at the protocol level—between application-layer defenses that fail at composition and infrastructure primitives that lack workflow semantics—enables uniform protection across heterogeneous frameworks while maintaining workflow context. All agent-resource interactions cross four fundamental control surfaces (prompts, tools, data, context); protecting these surfaces uniformly achieves compositional security impossible at either adjacent layer alone.

Three contributions realize this vision: MAPL provides cryptographic workflow dependencies via attestations, enabling sequential constraints without trusting application-reported state; embedded independent PEPs realize defense in depth where multiple cryptographic barriers must be defeated simultaneously; integration across nine frameworks through thin adapters validates surface completeness. Formal proofs (Lemmas 1-7, Theorems 1-3) establish security properties; empirical validation demonstrates 100\% recall, zero false positives across 174 test cases, and protection against two production CVEs that compromised baseline systems. This addresses the five challenges posed in Section~\ref{sec:introduction}: cryptographic integrity across frameworks, policy enforcement at scale, privilege non-escalation through composition, context integrity via hash chains, and complete accountability through non-repudiation. As agentic AI transitions to production, authenticated workflows provide the security substrate enabling safe deployment where composition gaps currently block adoption.
